% TC34 — คัดกรอง HIGH แจ้งที่ปรึกษา (ไม่สร้างคิว MH อัตโนมัติ)
\begin{testscenario}{34}{3}
  \tcmeta{TC34}{CS 2568/01}{ระบบเว็บแอปพลิเคชัน ``NUCARE''}{นายสมชาย ใจดี}{วันที่ 7 พฤศจิกายน พ.ศ.\ 2569}{1.2}
  \tcname{ผลคัดกรองความเสี่ยงสูงแจ้งที่ปรึกษาทางแอปและอีเมล}
  \tcpre{นิสิตมีที่ปรึกษาหรือไม่มีก็ได้ · รัน unit HRS-N · EMAIL\_PROVIDER พร้อมหรือ console}
  \tcdesc{เมื่อผลแบบประเมินเป็น HIGH ระบบแจ้งที่ปรึกษา (หรือ SD/ADMIN ถ้ายังไม่มอบหมาย) ไม่สร้าง CaseReferral ไป MH อัตโนมัติ}

  \tccase{1}{วางแผนแจ้งที่ปรึกษา}{%
    \begin{tcsteps}
      \item รัน \tccode{high-risk-screening-notify.test.ts}
    \end{tcsteps}}{%
    มีที่ปรึกษา → notify\_advisors · ไม่มี → notify\_staff\_unassigned}{%
    HRS-N01..02 ผ่าน}{Pass}

  \tccase{2}{หลังส่งแบบ HIGH ไม่โผล่คิว MH อัตโนมัติ}{%
    \begin{tcsteps}
      \item นิสิตส่งแบบประเมินได้ผล HIGH
      \item เข้าพอร์ทัล MH ดูคิวรอรับ
    \end{tcsteps}}{%
    ไม่มี referral จากระบบ · ที่ปรึกษาได้รับ alert/อีเมล}{%
    \tccode{notifyHighRiskScreening} ไม่เรียกสร้าง CaseReferral}{Pass}

  \tccase{3}{ไม่สร้าง EmergencyCase จากคัดกรอง}{%
    \begin{tcsteps}
      \item ตรวจคิวช่วยเหลือด่วน (SOS) หลังผล HIGH
    \end{tcsteps}}{%
    ไม่โผล่เป็น SOS ของนิสิต}{%
    screening ไม่สร้าง EmergencyCase}{Pass}
\end{testscenario}
